% Generated draft: requires mathematical and source review.
\documentclass[11pt]{article}
\usepackage[margin=25mm]{geometry}
\usepackage{fontspec}
\setmainfont{Latin Modern Roman}
\usepackage{amsmath,amssymb,mathtools,unicode-math}
\setmathfont{Latin Modern Math}
\usepackage{enumitem,needspace,xurl,hyperref}
\setlength{\parindent}{0pt}
\setlength{\parskip}{5pt}
\setlength{\emergencystretch}{3em}
\newcommand{\N}{\mathbb{N}}
\newcommand{\Z}{\mathbb{Z}}
\newcommand{\Q}{\mathbb{Q}}
\newcommand{\R}{\mathbb{R}}
\newcommand{\C}{\mathbb{C}}
\newcommand{\F}{\mathbb{F}}
\newcommand{\A}{\mathbb{A}}
\newcommand{\PP}{\mathbb{P}}
\newcommand{\E}{\mathbb{E}}
\newcommand{\eps}{\varepsilon}
\newcommand{\dd}{\,\mathrm{d}}
\begin{document}
\section*{UnsolvedMath ID: 30001694}
\textbf{Source problem code:} catalogue entry
\subsection{Definitions and mathematical statement}
A polyomino is understood here in the usual unit-grid sense: a nonempty finite union of closed squares \([i,i+1]\times[j,j+1]\), with \((i,j)\in\Z^2\), connected through shared edges. Let \(P\) be such a polyomino whose boundary \(J=\partial P\) is a Jordan curve, and let \(\Omega\) be the bounded component of \(\R^2\setminus J\). Define the greatest side length of an axis-parallel open square in \(\Omega\) by
\[
s=\max\bigl\{t>0:\text{there exist }x,y\in\R\text{ with }(x,x+t)\times(y,y+t)\subseteq\Omega\bigr\}.
\]
The supplied conjecture asks whether, for every such \(P\), there exist \(v,a,b\in\Z^2\) satisfying
\[
a\cdot b=0,\qquad \|a\|_2=\|b\|_2=\ell>0,
\qquad v,\ v+a,\ v+a+b,\ v+b\in J,
\]
with
\[
\ell\ge
\begin{cases}
s/\sqrt{2},&s\in 2\Z_{>0},\\[2pt]
\sqrt{m^2+(m+1)^2},&s=2m+1,\quad m\in\Z_{\ge0}.
\end{cases}
\]
Thus the four points must be the vertices of a nondegenerate square; its orientation is unrestricted. Only the vertices are required to lie on \(J\); the statement imposes no additional containment condition on its sides or interior.

The source's even/odd formulation presupposes that \(s\) is integral, but does not separately state or justify this property. The lattice is interpreted above as the standard square lattice associated with the unit cells. These conventions should be confirmed during mathematical review.
\subsection{Short English statement}
For every polyomino with a simple closed boundary, must that boundary contain four lattice points forming a square whose side satisfies the stated parity-dependent lower bound in terms of the largest axis-parallel open square inside the boundary? The source implicitly treats that largest side length as an integer.
\subsection{Sources}
\begin{itemize}
\item Supplied local registry record: permanent UnsolvedMath ID \(30001694\), source problem code \texttt{catalogue entry}, title ``Large Inscribed Lattice Squares in Polyomino Boundaries,'' in the Oberwolfach Reports Open Problems collection. The supplied proposed year is 2011; no individual proposer is identified.
\item Supplied source citation: \emph{Topological and Geometric Combinatorics} (2011), Oberwolfach Reports, pp.~349--423, DOI \texttt{10.4171/owr/2011/08}. Original OWR report: \url{https://doi.org/10.4171/owr/2011/08}.
\item Supplied literature reference, labeled ``grid-square computational study'': \url{https://doi.org/10.1007/s00454-014-9578-5}. No authors, article title, or precise computational bounds are supplied.
\item Literature assessment dated exactly \texttt{2026-08-21}: the supplied curation labels the problem open and says that the polyomino conjecture was computationally verified only through bounded grid sizes and remains unresolved. This is supplied dated curation, not an independently verified literature assessment. The background explicitly requests independent verification of time-sensitive status and cited claims.
\item The registry labels the statement \texttt{corrected\_verified}, describing a manual semantic, typographic, and LaTeX review. This supplied label does not constitute mathematical review of this draft or independent verification of the cited sources.
\end{itemize}
\end{document}
